\documentclass[11pt,a4paper]{article}
% preamble.tex — shared setup for human-verification documents
% Compile each document with:  tectonic <File>.tex

\usepackage[T1]{fontenc}
\usepackage[utf8]{inputenc}
\usepackage{lmodern}
\usepackage{amssymb,amsmath,amsthm}
\usepackage{mathtools}
\usepackage[dvipsnames]{xcolor}
\usepackage{listings}
\usepackage{hyperref}
\usepackage[margin=2.5cm]{geometry}
\usepackage{parskip}

% ── Colours ──────────────────────────────────────────────────────────
\definecolor{leanBg}{HTML}{F7F7F0}
\definecolor{haskellBg}{HTML}{F0F4F7}
\definecolor{keyword}{HTML}{7B3294}
\definecolor{comment}{HTML}{5A8C32}
\definecolor{noteFg}{HTML}{B35900}

% ── Lean 4 listings language ─────────────────────────────────────────
% NOTE: Lean listing blocks use ASCII approximations of unicode symbols
% (e.g. -> for arrow, o. for compose).  The literate key substitutes
% these with proper math-mode renderings in the PDF.
\lstdefinelanguage{Lean4}{
  morekeywords={def,theorem,structure,where,namespace,end,import,fun,
                let,have,show,by,exact,intro,cases,rfl,subst,obtain,
                funext,simp,rw,apply,variable,Type,Prop,instance,
                section,open,class,extends,deriving,sorry,noncomputable},
  sensitive=true,
  morecomment=[l]{--},
  morecomment=[n]{/-}{-/},
  morestring=[b]",
}

% ── Lean listing style ───────────────────────────────────────────────
\lstdefinestyle{lean}{
  language=Lean4,
  backgroundcolor=\color{leanBg},
  basicstyle=\small\ttfamily,
  keywordstyle=\color{keyword}\bfseries,
  commentstyle=\color{comment}\itshape,
  breaklines=true,
  frame=single,
  framerule=0.4pt,
  xleftmargin=4pt,
  xrightmargin=4pt,
  aboveskip=6pt,
  belowskip=4pt,
  literate=
    {->}{{$\to$}}2
    {=>}{{$\Rightarrow$}}2
    {<|}{{$\langle$}}1
    {|>}{{$\rangle$}}1
    {==>}{{$\Longrightarrow$}}3
    { o. }{{$\;\circ\;$}}3
    {forall}{{$\forall$}}1
    {:=}{{$\,:=\,$}}2
    {!=}{{$\neq$}}2
    {<=}{{$\leq$}}2
    {>=}{{$\geq$}}2,
}

% ── Haskell listing style ───────────────────────────────────────────
\lstdefinestyle{haskell}{
  language=Haskell,
  backgroundcolor=\color{haskellBg},
  basicstyle=\small\ttfamily,
  keywordstyle=\color{keyword}\bfseries,
  commentstyle=\color{comment}\itshape,
  breaklines=true,
  frame=single,
  framerule=0.4pt,
  xleftmargin=4pt,
  xrightmargin=4pt,
  aboveskip=6pt,
  belowskip=4pt,
  literate=
    {->}{{$\to$}}2
    {=>}{{$\Rightarrow$}}2,
}

% ── Theorem environments ────────────────────────────────────────────
\theoremstyle{definition}
\newtheorem{defn}{Definition}[section]
\newtheorem{thm}[defn]{Theorem}
\newtheorem{prop}[defn]{Proposition}

\theoremstyle{remark}
\newtheorem*{fidelitynote}{Fidelity note}

% ── Convenience macros ───────────────────────────────────────────────
\newcommand{\lean}[1]{\texttt{#1}}      % inline Lean identifier
\newcommand{\hask}[1]{\texttt{#1}}      % inline Haskell identifier
\newcommand{\Ext}{\mathrm{Ext}}
\newcommand{\Pos}{\mathrm{Pos}}
\newcommand{\Shape}{\mathrm{Shape}}


\title{Human Verification: \texttt{Containers.Directed}\\[4pt]
  \large Directed Containers and the Comonad They Induce}
\author{Machine-generated from \texttt{lean/Directed.lean} for review by Neil Ghani}
\date{July 2026}

\begin{document}
\maketitle

\begin{abstract}
This document presents every definition and theorem \emph{statement} from the
Lean~4 module \lean{Containers.Directed} in three parallel forms: the original
Lean, a readable Haskell rendering, and precise mathematical English.  Proofs
are omitted --- they are machine-checked by Lean.  The purpose is to let a
human verifier confirm that the formalisation captures the intended mathematics.
\end{abstract}

\tableofcontents
\bigskip

%===================================================================
\section{Overview}
%===================================================================

This file develops \emph{directed containers} (Ahman--Chapman--Uustalu,
\emph{When is a container a comonad?}) and shows that the extension $\Ext\,C$
of a directed container $C$ carries a canonical comonad structure.

A directed container augments a container $\langle S, P \rangle$ with:
\begin{itemize}
  \item a \emph{root} $o(s) \in P(s)$ for each shape $s$;
  \item a \emph{sub-shape} operation $s \downarrow p \in S$ for each position
        $p \in P(s)$;
  \item a \emph{shift} $p \oplus q \in P(s)$ combining $p \in P(s)$ with
        $q \in P(s \downarrow p)$;
  \item five laws D1--D5 making each $P(s)$ a ``monoid over the base''.
\end{itemize}

\noindent The comonad data on $D := \Ext\,C$ are:
\[
  \varepsilon\langle s, v \rangle = v(o(s)), \qquad
  \delta\langle s, v \rangle = \langle s, p \mapsto \langle s \downarrow p,
    q \mapsto v(p \oplus q) \rangle \rangle.
\]

The three comonad laws follow from D1--D5.  Laws D2 and D5 carry explicit
dependent transports (along D1 and D4 respectively) --- this is the genuinely
dependent content of the development.

%===================================================================
\section{The extensionality lemma}
%===================================================================

% --- 1. Ext.ext_eq --------------------------------------------------
\subsection{Theorem: \lean{Ext.ext\_eq}}

\begin{thm}[Extensionality for $\Ext$ across a shape equality]
Let $s_1, s_2 : S$ with $\lean{hs} : s_1 = s_2$, and let
$v_1 : P(s_1) \to X$, $v_2 : P(s_2) \to X$.  If the labellings agree after
transporting positions along \lean{hs},
\[
  \forall\, q : P(s_1),\quad v_1(q) = v_2(\lean{hs} \triangleright q),
\]
then $\langle s_1, v_1 \rangle = \langle s_2, v_2 \rangle$ as elements of
$\Ext\,C\,X$.
\end{thm}

\paragraph{Lean.}
\begin{lstlisting}[style=lean]
theorem Ext.ext_eq {C : Container} {X : Type} {s1 s2 : C.Shape}
    (hs : s1 = s2) {v1 : C.Pos s1 -> X} {v2 : C.Pos s2 -> X}
    (hv : forall q : C.Pos s1, v1 q = v2 (hs |-> q)) :
    (<|s1, v1|> : Ext C X) = <|s2, v2|>
\end{lstlisting}

\paragraph{Haskell.}
\begin{lstlisting}[style=haskell]
-- In Haskell, Ext s p x = (s, p -> x) and we have no shape equality
-- to transport along.  The analogue is simply:
-- if s1 == s2 and forall q, v1 q == v2 q, then (s1, v1) == (s2, v2).
\end{lstlisting}

\begin{fidelitynote}
The expression $\lean{hs} \triangleright q$ is Lean's \emph{transport}: it
coerces $q : P(s_1)$ to type $P(s_2)$ along the equality $s_1 = s_2$.  This
is the central dependent-type operation in the entire file.  It appears
whenever positions in different fibres $P(s)$ must be compared.  There is no
Haskell analogue; in Haskell all position types are the same (a single type
parameter \hask{p}), so transports are invisible.
\end{fidelitynote}


%===================================================================
\section{Directed containers}
%===================================================================

% --- 2. DirectedContainer --------------------------------------------
\subsection{Structure: \lean{DirectedContainer}}

\begin{defn}[Directed container]
A \textbf{directed container} extends a container $\langle S, P \rangle$ with:
\begin{itemize}
  \item $\mathrm{root} : (s : S) \to P(s)$ --- the root position of each shape;
  \item $\mathrm{sub} : (s : S) \to P(s) \to S$ --- the sub-shape
        $s \downarrow p$ beneath position $p$;
  \item $\mathrm{shift} : (s : S) \to (p : P(s)) \to P(s \downarrow p) \to P(s)$
        --- the shift $p \oplus q$;
\end{itemize}
satisfying the five laws:
\begin{itemize}
  \item[D1.] $s \downarrow o(s) = s$ \hfill (sub-shape under the root is the
             shape itself);
  \item[D2.] $\forall s, q : P(s \downarrow o(s)),\;
             o(s) \oplus q = \mathrm{D1}(s) \triangleright q$
             \hfill (shifting by the root is transport along D1);
  \item[D3.] $\forall s, p : P(s),\;
             p \oplus o(s \downarrow p) = p$
             \hfill (shift by sub-root is identity);
  \item[D4.] $\forall s, p, q,\;
             s \downarrow (p \oplus q) = (s \downarrow p) \downarrow q$
             \hfill (sub-shapes nest along shifts);
  \item[D5.] $\forall s, p, q, q',\;
             (p \oplus q) \oplus (\mathrm{D4}^{-1} \triangleright q')
             = p \oplus (q \oplus q')$
             \hfill (shift is associative, with transport along D4).
\end{itemize}
\end{defn}

\paragraph{Lean.}
\begin{lstlisting}[style=lean]
structure DirectedContainer extends Container where
  root  : (s : Shape) -> Pos s
  sub   : (s : Shape) -> Pos s -> Shape
  shift : (s : Shape) -> (p : Pos s) -> Pos (sub s p) -> Pos s
  d1 : forall s, sub s (root s) = s
  d4 : forall (s : Shape) (p : Pos s) (q : Pos (sub s p)),
    sub s (shift s p q) = sub (sub s p) q
  d2 : forall (s : Shape) (q : Pos (sub s (root s))),
    shift s (root s) q = d1 s |-> q
  d3 : forall (s : Shape) (p : Pos s),
    shift s p (root (sub s p)) = p
  d5 : forall (s : Shape) (p : Pos s) (q : Pos (sub s p))
              (q' : Pos (sub (sub s p) q)),
    shift s (shift s p q) ((d4 s p q).symm |-> q')
      = shift s p (shift (sub s p) q q')
\end{lstlisting}

\paragraph{Haskell.}
\begin{lstlisting}[style=haskell]
-- A directed container approximated in Haskell
-- (no dependent types; s and p are type parameters).
class (Container s p) => DirectedContainer s p where
  root  :: s -> p
  sub   :: s -> p -> s
  shift :: s -> p -> p -> p   -- shift s p q : p, not Pos (sub s p)
  -- Laws (as QuickCheck properties):
  -- D1: sub s (root s) == s
  -- D2: shift s (root s) q == q   (transport invisible in Haskell)
  -- D3: shift s p (root (sub s p)) == p
  -- D4: sub s (shift s p q) == sub (sub s p) q
  -- D5: shift s (shift s p q) q' == shift s p (shift (sub s p) q q')
\end{lstlisting}

\begin{fidelitynote}
In Lean, the shift has the precise dependent type
$\mathrm{shift}(s, p, -)  : P(s \downarrow p) \to P(s)$, and laws D2 and D5
carry transport coercions (along D1 and D4 respectively) that coerce positions
between fibres.  In the Haskell rendering, all positions are the same type
\hask{p}, so these transports are invisible --- the laws D2 and D5 still hold
but without the coercions.  The mathematical content (the monoid-over-the-base
structure on $P(s)$) is preserved; the dependent bookkeeping is not.
\end{fidelitynote}


%===================================================================
\section{The comonad on \lean{Ext C}}
%===================================================================

% --- 3. DirectedContainer.counit ------------------------------------
\subsection{Definition: \lean{DirectedContainer.counit}}

\begin{defn}[Counit]
The \textbf{counit} $\varepsilon : \Ext\,C \Rightarrow \mathrm{Id}$ at $X$ is:
\[
  \varepsilon\langle s, v \rangle \;=\; v(o(s)).
\]
It evaluates a labelled container at the root position of its shape.
\end{defn}

\paragraph{Lean.}
\begin{lstlisting}[style=lean]
def DirectedContainer.counit : Ext D.toContainer X -> X :=
  fun p => p.2 (D.root p.1)
\end{lstlisting}

\paragraph{Haskell.}
\begin{lstlisting}[style=haskell]
counit :: DirectedContainer s p => (s, p -> x) -> x
counit (s, v) = v (root s)

-- As a comonad extract:
instance DirectedContainer s p => Comonad (Ext s p) where
  extract (Ext s v) = v (root s)
\end{lstlisting}


% --- 4. DirectedContainer.comult ------------------------------------
\subsection{Definition: \lean{DirectedContainer.comult}}

\begin{defn}[Comultiplication]
The \textbf{comultiplication} $\delta : \Ext\,C \Rightarrow \Ext\,C \circ \Ext\,C$
at $X$ is:
\[
  \delta\langle s, v \rangle \;=\;
  \Bigl\langle s,\; p \mapsto \bigl\langle s \downarrow p,\;
    q \mapsto v(p \oplus q) \bigr\rangle \Bigr\rangle.
\]
Each position $p$ yields a new labelled container of sub-shape $s \downarrow p$,
whose labelling evaluates $v$ at the shift $p \oplus q$.
\end{defn}

\paragraph{Lean.}
\begin{lstlisting}[style=lean]
def DirectedContainer.comult :
    Ext D.toContainer X -> Ext D.toContainer (Ext D.toContainer X) :=
  fun p => <|p.1, fun q => <|D.sub p.1 q, fun r => p.2 (D.shift p.1 q r)|>|>
\end{lstlisting}

\paragraph{Haskell.}
\begin{lstlisting}[style=haskell]
comult :: DirectedContainer s p => (s, p -> x) -> (s, p -> (s, p -> x))
comult (s, v) = (s, \q -> (sub s q, \r -> v (shift s q r)))

-- As a comonad duplicate:
  duplicate (Ext s v) = Ext s (\q -> Ext (sub s q) (\r -> v (shift s q r)))
\end{lstlisting}


% --- 5. DirectedContainer.left_counit --------------------------------
\subsection{Theorem: \lean{DirectedContainer.left\_counit}}

\begin{thm}[Left counit law]
\[
  \varepsilon \circ \delta \;=\; \mathrm{id}_{\Ext\,C\,X}.
\]
The shape equation uses D1; the value equation uses D2 (with a transport).
\end{thm}

\paragraph{Lean.}
\begin{lstlisting}[style=lean]
theorem DirectedContainer.left_counit :
    D.counit o. D.comult
      = (id : Ext D.toContainer X -> Ext D.toContainer X)
\end{lstlisting}

\paragraph{Haskell.}
\begin{lstlisting}[style=haskell]
-- Comonad law (left unit):
-- extract . duplicate = id
--
-- Proof: (extract . duplicate) (s, v)
--   = extract (s, \q -> (sub s q, \r -> v (shift s q r)))
--   = let (sub s (root s), \r -> v (shift s (root s) r))
--       = (s, \r -> v r)    -- by D1 and D2
--       = (s, v)
\end{lstlisting}


% --- 6. DirectedContainer.right_counit --------------------------------
\subsection{Theorem: \lean{DirectedContainer.right\_counit}}

\begin{thm}[Right counit law]
\[
  (\Ext\,C\,\varepsilon) \circ \delta \;=\; \mathrm{id}_{\Ext\,C\,X}.
\]
The shapes are unchanged; the value equation uses D3.
\end{thm}

\paragraph{Lean.}
\begin{lstlisting}[style=lean]
theorem DirectedContainer.right_counit :
    Ext.map D.counit o. D.comult
      = (id : Ext D.toContainer X -> Ext D.toContainer X)
\end{lstlisting}

\paragraph{Haskell.}
\begin{lstlisting}[style=haskell]
-- Comonad law (right unit):
-- fmap extract . duplicate = id
--
-- Proof: (fmap extract . duplicate) (s, v)
--   = fmap extract (s, \q -> (sub s q, \r -> v (shift s q r)))
--   = (s, \q -> extract (sub s q, \r -> v (shift s q r)))
--   = (s, \q -> v (shift s q (root (sub s q))))
--   = (s, \q -> v q)    -- by D3
--   = (s, v)
\end{lstlisting}


% --- 7. DirectedContainer.coassoc ------------------------------------
\subsection{Theorem: \lean{DirectedContainer.coassoc}}

\begin{thm}[Coassociativity]
\[
  (\Ext\,C\,\delta) \circ \delta
  \;=\;
  \delta \circ \delta
  \;:\;
  \Ext\,C\,X \to \Ext\,C\,(\Ext\,C\,(\Ext\,C\,X)).
\]
The inner shapes agree by D4; the deepest values agree by D5, transported along
D4.
\end{thm}

\paragraph{Lean.}
\begin{lstlisting}[style=lean]
theorem DirectedContainer.coassoc :
    Ext.map D.comult o. D.comult
      = (D.comult o. D.comult :
          Ext D.toContainer X
            -> Ext D.toContainer (Ext D.toContainer (Ext D.toContainer X)))
\end{lstlisting}

\paragraph{Haskell.}
\begin{lstlisting}[style=haskell]
-- Comonad law (coassociativity):
-- fmap duplicate . duplicate = duplicate . duplicate
--
-- Proof: both sides send (s, v) to
--   (s, \p -> (sub s p, \q -> (sub (sub s p) q, \r -> v (shift s p (shift (sub s p) q r)))))
-- using D4 (sub s (shift s p q) = sub (sub s p) q) and
--      D5 (shift s (shift s p q) r = shift s p (shift (sub s p) q r)).
-- In Haskell these transports are invisible; in Lean they require ext_eq.
\end{lstlisting}

\begin{fidelitynote}
The coassociativity proof in Lean requires two applications of \lean{Ext.ext\_eq}
to deal with the shape equality D4 and the transport in D5.  Specifically:
the left-hand side has inner shape $s \downarrow (p \oplus q)$, while the
right-hand side has inner shape $(s \downarrow p) \downarrow q$; these are
equal by D4 but not definitionally.  D5 then equates the value functions after
transporting the position $q'$ along D4.  In the Haskell rendering (where all
positions are a single type), these transports vanish and the law holds by
simple equational reasoning on $\mathrm{shift}$.
\end{fidelitynote}


%===================================================================
\section{Translation conventions}
%===================================================================

\begin{enumerate}
\item \textbf{Transport ($\triangleright$).}  The Lean notation \lean{h} $\triangleright$ \lean{q}
  transports $q$ along equality $h$.  This appears in D2, D5, and the proof
  of \lean{Ext.ext\_eq}.  It is the most important dependent-type feature of
  this file and has no Haskell counterpart.

\item \textbf{Sub-shape types.}  In Lean, $q : P(s \downarrow p)$ --- the type
  of $q$ depends on both $s$ and $p$.  In Haskell we approximate this by a
  single position type \hask{p}, losing the fibre structure.

\item \textbf{Comonad laws.}  The three comonad laws are stated as equalities
  of functions $\Ext\,C\,X \to -$.  In Haskell, these are the standard
  comonad laws \hask{extract . duplicate = id},
  \hask{fmap extract . duplicate = id}, and
  \hask{fmap duplicate . duplicate = duplicate . duplicate}.

\item \textbf{Universe level.}  The Lean development is entirely at
  $\mathbf{Type}$ (a single universe) with no Mathlib dependency.
\end{enumerate}

\end{document}
